\chapter{Guidelines and Grading Criteria for Exercises}
\label{app:petunjuk}

The answer key is in Appendix~D, but its form is deliberately incomplete: final numbers
and sketches of the steps, not full proofs. The reason is consistency. A final number
without a trace of checking will teach the reader the very habit this book wants to
remove, and a proof served up whole removes the only useful exercise, namely finding for
yourself where an argument does not yet close. So use this key as a last comparison, not
as a place to start. What is provided here is a way to check your own work before you
open Appendix~D.

\section*{Three kinds of exercises and how to check them}

\begin{center}
\small
\begin{tabular}{@{}p{2.6cm}p{5.6cm}p{5.8cm}@{}}
\toprule
\textbf{Type} & \textbf{What is assessed} & \textbf{How to check} \\
\midrule
Proof & every step must start from a definition or a result already
  proved; no leaps are allowed & swap drafts with a classmate; ask her
  to point to the sentence she cannot derive herself \\
Computation & the final number and the intermediate steps & redo it with Python; compute the relative
  error against the exact result; check whether the answer makes sense in
  order of magnitude \\
Evaluating a claim & the quality of the argument, not a binary true-false & first write down what
  must be true for the claim to hold; look for one example that breaks it;
  only then decide ``true'', ``false'', or ``not enough evidence'' \\
\bottomrule
\end{tabular}
\end{center}

\section*{Verification tools per part}

\begin{center}
\small
\begin{tabular}{@{}llp{7.2cm}@{}}
\toprule
\textbf{Chapter} & \textbf{Tool} & \textbf{Typical check} \\
\midrule
1--5 & \texttt{python3} & table of values, factorization, partial sums; compare
  with an analytic estimate (e.g. $H_n\approx\ln n+\gamma$) \\
6--10 & \texttt{numpy}, \texttt{scipy.linalg} & rank, null space (\texttt{null\_space}),
  eigenvalues, SVD, least squares (\texttt{lstsq}); check
  $\|A\mathbf{v}-\lambda\mathbf{v}\|$ \\
11--16 & \texttt{scipy.integrate}, finite differences & quadrature vs antiderivative; rate of
  error at several $h$ or several $n$ \\
17--21 & \texttt{numpy.random}, \texttt{scipy.stats} & Monte Carlo simulation with a
  fixed \emph{seed}; compare the empirical probability with the exact probability and
  report its standard error \\
22 & \texttt{struct}, \texttt{float.hex} & bitfields, ULP, machine epsilon;
  compare the naive algorithm with Kahan \\
23--24 & \texttt{scipy.optimize} & \texttt{linprog}, \texttt{minimize},
  \texttt{milp}; check the KKT conditions and strong duality numerically \\
25--26 & \texttt{numpy} & gradient check by finite differences; training vs validation
  error curves at several regularization settings \\
27--30 & \texttt{numpy}, \texttt{scikit-learn} when available & confusion matrix from
  stored predictions; Wilson interval for proportions; accuracy under attack
  for several $\varepsilon$ and several \emph{seed} \\
\bottomrule
\end{tabular}
\end{center}

\section*{How to grade your own work}

\begin{enumerate}
  \item First write down what is given and what is to be proved.
        If neither is clear, the difficulty is not in the mathematics.
  \item For computation exercises, guess the answer before computing. A wrong guess
        is still useful: it shows where intuition misses.
  \item Compare the result by two different methods. One method can be right by
        accident; two methods that agree are rarely an accident.
  \item If the results differ greatly, look for the mistake in a simple case first
        (one variable, one step, $n=1$).
  \item Write down the limits of validity of the answer. An answer that is correct for a
        special case but presented as a general answer is the mistake that most often slips through.
\end{enumerate}
